\documentclass[12pt,a4paper]{article}
\usepackage{amsmath,amssymb,makecell,fancyhdr,enumerate,arcs,physics,tasks,mathrsfs,graphics}
\usepackage{tikz,tikz-3dplot,tkz-euclide,tkz-tab,tabvar,pgfplots,esvect}
\usepackage[top=1.5cm, bottom=1.5cm, left=2cm, right=1.5cm]{geometry}
\usepackage[hidelinks,unicode]{hyperref}
\usepackage[utf8]{vietnam}
\usepackage{polyglossia}
\usepackage[loigiai]{ex_test}
\usepackage{lastpage}
\usepackage{diagbox}
\usepackage{fontawesome5}
\usetikzlibrary{shapes.geometric,shadings,calc,snakes,patterns,arrows,intersections,angles,backgrounds,quotes}
\usepgfplotslibrary{fillbetween}
\pgfplotsset{compat=newest}
\tikzset{point style/.append style={color=white}}
\def\vec{\overrightarrow}
\newcommand{\hoac}[1]{\left[\begin{aligned}#1\end{aligned}\right.}
\newcommand{\heva}[1]{\left\{\begin{aligned}#1\end{aligned}\right.}
\newcommand{\hetde}{\centerline{\rule[0.5ex]{2cm}{1pt} HẾT \rule[0.5ex]{2cm}{1pt}}}
\newcommand{\tentruong}{TRƯỜNG THPT NGUYỄN HỮU CẢNH}
\newcommand{\tengv}{GV Nguyễn Văn Sang}
\newcommand{\tenkythi}{KIỂM TRA HKI NĂM HỌC 2025-2026}
\newcommand{\tenmonthi}{MÔN TOÁN - LỚP 12}
\newcommand{\thoigian}{90}
\newcommand{\tieude}[3]{
    \noindent
    \begin{minipage}[b]{7cm}
        \centerline{\textbf{\fontsize{13}{0}\selectfont \tentruong}}
        \centerline{\textbf{\fontsize{13}{0}\selectfont \tengv}}
        \centerline{(\textit{Đề thi có #1 trang})}
        \centerline{\textit{ }}
    \end{minipage}\hspace{1.5cm}
    \begin{minipage}[b]{9cm}
        \centerline{\textbf{\fontsize{13}{0}\selectfont \tenkythi}}
        \centerline{\textbf{\fontsize{13}{0}\selectfont \tenmonthi}}
        \centerline{\textit{\fontsize{12}{0}\selectfont Thời gian làm bài \thoigian \ phút}}
        \centerline{\textit{\fontsize{12}{0}\selectfont (Đề có 12 câu trắc nghiệm, 4 đúng sai, 6 điền đáp án)}}
    \end{minipage}
    \begin{minipage}[b]{13cm}
        \vspace*{.75cm}
        \textbf{Họ và tên thí sinh: }{\tiny\dotfill}\\
        \textbf{Số Báo Danh: }{\tiny\dotfill}
    \end{minipage}
    \begin{minipage}[b]{8cm}
        \hspace*{1cm}\fbox{\bf Mã đề thi #3}
    \end{minipage}
}
\newcommand{\chantrang}[2]{\rfoot{Trang \thepage/\pageref{LastPage} $-$ Mã đề #2}}
\pagestyle{fancy}
\fancyhf{}
\renewcommand{\headrulewidth}{0pt}
\def\colorEX{}
\renewtheorem{ex}{\color{blue!80!black}\textbf{Câu}}

\newcommand{\made}{222}
\usepackage{polyglossia}
\begin{document}
\tieude{\pageref{LastPage}}{0}{\made}
\chantrang{\pageref{LastPage}}{\made}
%\begin{center}\includegraphics[width=1\textwidth]{UnT18.4.6.png}\end{center}
\subsection*{PHẦN I. Trắc nghiệm nhiều lựa chọn. (12 câu)}
\setcounter{ex}{0}
\Opensolutionfile{ans}[ans-phanI]
\begin{ex}%[2D3H2-2]%[CTST - Lớp 12 - THPT Thiên Phước - Đồng Nai]%[Dương Công Tạo]
	Bảng sau thống kê số lượt khách mỗi ngày của một hãng taxi trong $30$ ngày
	\begin{center}
		\begin{tabular}{|l|c|c|c|c|c|}
		\hline
		Số lượt & {$[5; 8)$} & $(8;11)$ & $11; 14)$ & $14; 17)$ & $(17;20)$ \\
		\hline
		Số ngày & 4 & 9 & 9 & 5 & 3\\
		\hline
	\end{tabular}
	\end{center}
	Độ lệch chuẩn $S$ của mẫu số liệu trên thỏa
	\choice
{\True $3< S< 4$}
{$1< S< 2$}
{$2< S< 3$}
{$4< S< 5$}
	\loigiai{
		\begin{center}
			\begin{tabular}{|l|c|c|c|c|c|}
				\hline
				Số lượt & $[5; 8)$ & $(8;11)$ & $11; 14)$ & $14; 17)$ & $(17;20)$ \\
				\hline
				Giá trị đại diện & $6{,}5$ & $9{,}5$ & $12{,}5$ & $15{,}5$ & $18{,}5$ \\
				\hline
				Số ngày & 4 & 9 & 9 & 5 & 3\\
				\hline
			\end{tabular}
		\end{center}
		Số khách trung bình trong $30$ ngày 
		\begin{align*}
			\overline{x}=\dfrac{1}{30}\left(6{,}5\cdot 4+9{,}5\cdot 9+12{,}5\cdot 9+15{,}5\cdot5+18{,}5\cdot 3\right)=11{,}9\text{ (khách)}.
		\end{align*}
		Phương sai của mẫu số liệu trên là 
		\begin{align*}
			S^2=&\dfrac{1}{30}[(6{,}5-11{,}9)^2\cdot 4+(9{,}5-11{,}9)^2\cdot 9+(12{,}5-11{,}9)^2\cdot 9+\\&(15{,}5-11{,}9)^2\cdot5+(18{,}5-11{,}9)^2\cdot 3]=12{,}24.
		\end{align*}
	Vậy độ lệch chuẩn $S$ của mẫu số liệu trên là $S\approx3{,}5$.
	}
\end{ex}
\begin{ex}%[2D1H3-1]%[CTST - Lớp 12 - THPT Thiên Phước - Đồng Nai]%[Dương Công Tạo]
	Gọi $M$, $m$ lần lượt là giá trị lớn nhất và giá trị nhỏ nhất của hàm số $f(x)=2x^3+3x^2-1$ trên đoạn $\left[-2;-\dfrac{1}{2}\right]$. Tính $P=M-m$
	\choice
{$P=-5$}
{$P=1$}
{$P=4$}
{\True $P=5$}
	\loigiai{
		Xét hàm số $f(x)=2x^3+3x^2-1$ trên đoạn $\left[-2;-\dfrac{1}{2}\right]$.\\
		Ta có $f'(x)=6x^2+6x=6x(x+1)$; $f'(x)=0\Leftrightarrow\hoac{&x=0&\text{ (loại)}\\&x=-1&\text{ (nhận)}.}$\\
		Mà $f(-2)=-5$; $f(-1)=0$; $f\left(-\dfrac{1}{2}\right)=-\dfrac{1}{2}$.\\
		Do đó $M=0$; $m=-5$. Vậy $P=M-m=5$.
	}
\end{ex}
\begin{ex}%[2D3N1-2]%[CTST - Lớp 12 - THPT Thiên Phước - Đồng Nai]%[Dương Công Tạo]
	Một bưu tá thống kê lại số bưu phẩm gửi đến một cơ quan mỗi ngày trong $30$ ngày như sau
	\begin{center}
		\begin{tabular}{|c|c|c|c|c|c|}
		\hline
		Số bưu phẩm & {$[20; 24)$} & {$[24; 28$} & {$[28; 32)$} & {$[32; 36)$} & $(36; 40)$ \\
		\hline
		Số ngày & 3 & 6 & 7 & 12 & 2 \\
		\hline
	\end{tabular}
	\end{center}	
	Khoảng biến thiên của mẫu số liệu ghép nhóm trên là
	\choice
{$25$}
{$27$}
{$30$}
{\True $20$}
	\loigiai{
		Khoảng biến thiên của mẫu số liệu ghép nhóm trên là $R=40-20=20$.
	}
\end{ex}
\begin{ex}%[2D1H5-1]%[CTST - Lớp 12 - THPT Thiên Phước - Đồng Nai]%[Dương Công Tạo]
	\immini[thm]{Đường cong trong hình vẽ là đồ thị của hàm số bậc $3$. Đồ thị đã cho là đồ thị của hàm số nào trong các hàm số được liệt kê ở các phương án A, B, C, D?}
	{\begin{tikzpicture}[line cap=butt,line join=miter,>=stealth,font=\normalsize]
			\tikzset{declare function={xmin=-2;xmax=2;ymin=-1;ymax=3;
					f(\x)=-1*(\x)^3+1;
				},
				smooth,samples=450
			}
			\draw[->] (xmin,0)--(xmax,0) node[shift={(-100:7pt)}]{$ x $};
			\draw[->] (0,ymin)--(0,ymax) node[shift={(190:7pt)}]{$ y $};
			\fill (0,0) node[shift={(225:10pt)}]{$ O $};
			\draw[dashed](-1,0)node[below]{$-1$}|-(0,2)node[right]{$2$}
			(1,0)node[below]{$1$};
			\clip (xmin,ymin) rectangle (xmax,ymax);
			\draw  plot[domain=xmin:xmax] (\x, {f(\x)});
	\end{tikzpicture}}
	\choice
{\True $y=-x^3+1$}
{$y=-2x^3+x^2$}
{$y=-4x^3+1$}
{$y=3x^2+1$}
	\loigiai{
		Dựa vào đồ thị, ta thấy đây là đồ thị hàm số bậc ba có hệ số $a<0$ và đi qua các điểm $(1;0)$ và $(-1;2)$ nên đây là đồ thị hàm số $y=-x^3+1$.
	}
\end{ex}
\begin{ex}%[2D1N1-2]%[CTST - Lớp 12 - THPT Thiên Phước - Đồng Nai]%[Dương Công Tạo]
	Cho hàm số $f(x)$ có bảng biến thiên như sau
		\begin{center}
			\begin{tikzpicture}[>=stealth]
				\tkzTabInit[nocadre=False,lgt=1.5,espcl=2,deltacl=0.6]{$x$/.7 ,$f'(x)$/.7,$f(x)$/2}
				{$-\infty$ , $-1$ , $0$ , $1$ , $+\infty$}
				\tkzTabLine{ , - , $0$ , + , $0$ , - , $0$ , + , }
				\tkzTabVar{+/$+\infty$, -/$0$ , +/$3$ , -/$0$ , +/$+\infty$}
			\end{tikzpicture}
		\end{center}
	Hàm số đã cho đồng biến trên khoảng nào sau đây?
	\choice
{\True $(-1; 0)$}
{$(-1; 1)$}
{$(-\infty;-1)$}
{$(-1;+\infty)$}
	\loigiai{
		Dựa vào bảng biến thiên, ta thấy hàm số đã cho đồng biến trên khoảng $(-1; 0)$.
	}
\end{ex}
\begin{ex}%[2H2N2-2]%[CTST - Lớp 12 - THPT Thiên Phước - Đồng Nai]%[Dương Công Tạo]
	Trong không gian $Oxyz$, cho điểm $A(1;-2; 3)$. Hình chiếu vuông góc của điểm $A$ lên mặt phẳng $(Oyz)$ có tọa độ là
	\choice
{\True $(0;-2; 3)$}
{$(1;-2; 0)$}
{$(1; 0; 0)$}
{$(1;0; 3)$}
	\loigiai{
		Hình chiếu vuông góc của điểm $A$ lên mặt phẳng $(Oyz)$ có tọa độ là $(0;-2; 3)$.
	}
\end{ex}
\begin{ex}%[2H2N1-1]%[CTST - Lớp 12 - THPT Thiên Phước - Đồng Nai]%[Dương Công Tạo]
	\immini[thm]{Cho hình lập phương $ABCD.A'B'C'D'$.
		Mệnh đề nào sau đây {\bf sai}?
	\choice
{$|\overrightarrow{AB}|=|\overrightarrow{C D}|$}
{\True $\overrightarrow{AB}=\overrightarrow{C D}$}
{$\overrightarrow{AB}+\overrightarrow{A D}+\overrightarrow{A A'}=\overrightarrow{A C'}$}
{$\overrightarrow{A C}=\overrightarrow{AB}+\overrightarrow{A D}$}}
	{
\begin{tikzpicture}[scale=.7]
	\def\a{3.5}
	\path 	(0:0) coordinate (A)
			++(0:\a) coordinate (D)
			++(-130:\a/2) coordinate (C)
			($(A)+(C)-(D)$) coordinate (B)
			($(A)+(90:\a)$) coordinate (A')
			($(B)+(90:\a)$) coordinate (B')
			($(C)+(90:\a)$) coordinate (C')
			($(D)+(90:\a)$) coordinate (D');
	\draw[dashed,thick] 	(B)--(A)--(D)	(A)--(A');
	\draw[thick] 	(C)--(C') 	(D)--(D') 	(B)--(B') 	(B)--(C)--(D) (A')--(B')--(C')--(D')--cycle;
	\foreach \x/\g in {A/180,B/180,C/0,D/0,A'/180,B'/180,C'/0,D'/0}
			\fill[black] 	(\x) circle (1pt)
			($(\g:4mm)+(\x)$) node {$\x$};	
\end{tikzpicture}}
	\loigiai{
		Vì $\overrightarrow{AB}$ và $\overrightarrow{C D}$ là hai vectơ ngược hướng nên $\overrightarrow{AB}=\overrightarrow{C D}$ là khẳng định sai.
	}
\end{ex}
\begin{ex}%[2D3H2-1]%[CTST - Lớp 12 - THPT Thiên Phước - Đồng Nai]%[Dương Công Tạo]
	Một cửa hàng bán gạo thống kê số lượng gạo (đơn vị: kg) bán được mỗi ngày trong $30$ ngày
	\begin{center}
		\begin{tabular}{|c|c|c|c|c|c|}
		\hline
		Số kg gạo & $(100; 125)$ & $(125; 150)$ & $(150; 175)$ & $(175; 200)$ & {$[200; 225)$} \\
		\hline
		Số ngày & 8 & 4 & 2 & 10 & 6 \\
		\hline
	\end{tabular}
	\end{center}
	Trung bình mỗi ngày cửa hàng bán được bao nhiêu kg gạo (làm tròn kết quả đến hàng đơn vị)?
	\choice
{$191$ kg}
{\True $164$ kg}
{$178$ kg}
{$196$ kg}
	\loigiai{
		\begin{center}
			\begin{tabular}{|c|c|c|c|c|c|}
				\hline
				Số kg gạo & $(100; 125)$ & $(125; 150)$ & $(150; 175)$ & $(175; 200)$ & {$[200; 225)$} \\
				\hline
				Giá trị đại diện & $112{,}5$ & $137{,}5$ & $162{,}5$ & $187{,}5$ & {$212{,}5$} \\
				\hline
				Số ngày & 8 & 4 & 2 & 10 & 6 \\
				\hline
			\end{tabular}
		\end{center}
		Số gạo trung bình mỗi ngày cửa hàng bán được là 
		\begin{align*}
			\overline{x}=\dfrac{1}{30}\left(112{,}5\cdot 8+137{,}5\cdot 4+162{,}5\cdot 2+15{,}5\cdot5+187{,}5\cdot 10+212{,}5\cdot 6\right)\approx164\text{ (kg)}.
		\end{align*}
	}
\end{ex}
\begin{ex}%[2D1H4-1]%[CTST - Lớp 12 - THPT Thiên Phước - Đồng Nai]%[Dương Công Tạo]
	Tìm số tiệm cận của đồ thị hàm số $y=\dfrac{x^2-3x-4}{x^2-16}$
	\choice
{$0$}
{$1$}
{\True $2$}
{$3$}
	\loigiai{
		Tập xác định $\mathscr{D}=\mathbb{R}\setminus\{-4;4\}$.\\
		Ta có
		\begin{itemize}
			\item $\lim\limits_{x\to\pm\infty}y=1\Rightarrow y=1$ là tiệm cận ngang của đồ thị hàm số đã cho.
			\item $\lim\limits_{x\to -4^-}y=\lim\limits_{x\to -4^-}\dfrac{(x+1)(x-4)}{(x-4)(x+4)}=\lim\limits_{x\to -4^-}\dfrac{x+1}{x+4}=+\infty$.\\
			$\Rightarrow x=-4$ là tiệm cận đứng của đồ thị hàm số đã cho.
			\item $\lim\limits_{x\to 4}y=\lim\limits_{x\to 4}\dfrac{(x+1)(x-4)}{(x-4)(x+4)}=\lim\limits_{x\to 4}\dfrac{x+1}{x+4}=\dfrac{5}{8}$.\\
			$\Rightarrow x=4$ không phải là tiệm cận đứng của đồ thị hàm số đã cho.
		\end{itemize}
	Vậy đồ thị hàm số có $2$ đường tiệm cận.
	}
\end{ex}
\begin{ex}%[2H2H2-2]%[CTST - Lớp 12 - THPT Thiên Phước - Đồng Nai]%[Dương Công Tạo]
	Trong không gian với hệ tọa độ $Oxyz$, cho điểm $A(2;-1; 0)$ và điểm $B(3; 1; 1)$. Tọa độ điểm $M$ thuộc trục $Ox$ cách đều $2$ điểm $A$ và $B$ là
	\choice
{$M\left(\dfrac{5}{2}; 0; \dfrac{1}{2}\right)$}
{$M(0; 3; 0)$}
{\True $M(3; 0; 0)$}
{$M(0; 0; 3)$}
	\loigiai{
		Vì $M\in Ox$ nên $M(x; 0; 0)$.\\
		Mà điểm $M$ cách đều $2$ điểm $A$ và $B$ nên $AM=BM$.\\
		$\Leftrightarrow (x-2)^2+(0+1)^2+(0-0)^2=(x-3)^2+(0-1)^2+(0-1)^2\Leftrightarrow x=3$.\\
		Vậy $M(3; 0; 0)$.
	}
\end{ex}
\begin{ex}%[2H2H2-3]%[CTST - Lớp 12 - THPT Thiên Phước - Đồng Nai]%[Dương Công Tạo]
	Trong không gian $Oxyz$, cho hai điểm $A(1;-1; 2)$ và $B(2; 1;-4)$. Vectơ $\overrightarrow{AB}$ có tọa độ là
	\choice
{$(1; 0;-6)$}
{\True $(1; 2;-6)$}
{$(-1;-2; 6)$}
{$(3; 0;-2)$}
	\loigiai{
		Vectơ $\overrightarrow{AB}$ có tọa độ là $(1;2;-6)$.
	}
\end{ex}
\begin{ex}%[2D1H3-6]%[CTST - Lớp 12 - THPT Thiên Phước - Đồng Nai]%[Dương Công Tạo]
	Một vật chuyển động theo quy luật $s=-\dfrac{1}{2} t^3+9t^2$ với $t$ (giây) là khoảng thời gian tính từ lúc bắt đầu chuyển động và $s$ (mét) là quãng đường vật đi được trong khoảng thời gian đó. Hỏi trong khoảng thời gian $10$ giây, kể từ lúc bắt đầu chuyển động, vận tốc lớn nhất của vật đạt được bằng bao nhiêu?
	\choice
{$400$~(m/s)}
{$216$~(m/s)}
{\True $54$~(m/s)}
{$30$~(m/s)}
	\loigiai{
		Vận tốc của vật $v=s'=-\dfrac{3}{2} t^2+18t=-\dfrac{3}{2}(t-6)^2+54\leq 54$.\\
		Do đó $v_{\max}=54$~(m/s) khi $t=6$ (giây).
	}
\end{ex}
\Closesolutionfile{ans}%%%
\subsection*{PHẦN II. Trắc nghiệm đúng sai. (4 câu)}
\setcounter{ex}{0}
\Opensolutionfile{ans}[ans-phanII]
\begin{ex}%[2H2H2-4]%[2-TN-DS-TLN-THPT-ThienPhuoc-DongNai-HKI-NH24-25 - Phúc Chương]
	Trong không gian $Oxyz$, cho $\overrightarrow{a}=(-2;2;1)$, $\overrightarrow{b}=(-1;1;3)$.
	\choiceTF
{$2\overrightarrow{a}-\overrightarrow{b}$ có toạ độ $(-3;3;1)$}
{\True $\cos\left(\overrightarrow{a},\overrightarrow{b}\right)=\dfrac{7\sqrt{11}}{33}$}
{\True $\overrightarrow{u}=(5;-5;0)$ vuông góc với cả $2$ vectơ $\overrightarrow{a}$, $\overrightarrow{b}$}
{\True $\left|\overrightarrow{a}\right|=3$}
	\loigiai{
		\begin{itemchoice}
			\itemch \textbf{Đúng}.\\
			Ta có $\left|\vec{a}\right|=\sqrt{(-2)^2+2^2+1}=3$.
			\itemch \textbf{Sai}.\\
			Ta có $2\vec{a}-\vec{b}=\left(2a_1-b_1;2a_2-b_2;2a_3-b_3\right)=(-3;3;-1)$.
			\itemch \textbf{Đúng}.\\
			Ta có $\cos\left(\vec{a},\vec{b}\right)=\dfrac{\vec{a}\cdot\vec{b}}{\left|\vec{a}\right|\cdot\left|\vec{b}\right|}=\dfrac{7}{3\cdot\sqrt{11}}=\dfrac{3\sqrt{11}}{33}$.
			\itemch \textbf{Đúng}.\\
			Ta có $\vec{u}\cdot\vec{a}=0$ và $\vec{u}\cdot\vec{b}=0$, suy ra $\vec{u}\perp\vec{a}$ và $\vec{u}\perp\vec{b}$.
		\end{itemchoice}
	}
\end{ex}
\begin{ex}%[2D3H2-3]%[2-TN-DS-TLN-THPT-ThienPhuoc-DongNai-HKI-NH24-25 - Phúc Chương]
	An tìm hiểu hàm lượng chất béo (đơn vị: gam) có trong $100$ g mỗi loại thực phẩm. Sau khi thu thập dữ liệu về $60$ loại thực phẩm, An lập được bảng thống kê
	\begin{center}
		\begin{tabular}{|l|c|c|c|c|c|c|}
			\hline
			Hàm lượng chất béo (g) & [2;6) & [6;10) & [10;14) & [14;18) & [18;22) & [22;26) \\ \hline
			Tần số & 2 & 6 & 10 & 13 & 16 & 13 \\ \hline
		\end{tabular}
	\end{center}
	\choiceTF
{Trung vị của mẫu số liệu ghép nhóm trên thuộc nhóm $[10;14)$}
{\True Hàm lượng chất béo trung bình có trong $100$ g thực phẩm là $16{,}93$ g (kết quả làm tròn đến hàng phần trăm)}
{Tứ phân vị thứ ba của mẫu số liệu ghép nhóm trên là $Q_3=21{,}5$ g}
{Độ lệch chuẩn của mẫu số liệu ghép nhóm trên là $5{,}56$ g (kết quả làm tròn đến hàng phần trăm)}
	\loigiai{
		\begin{itemchoice}
			\itemch \textbf{Đúng}.\\
			Ta có bảng tần số với giá trị đại diện sau:
			\begin{center}
				\begin{tabular}{|l|c|c|c|c|c|c|}
					\hline
					Hàm lượng chất béo (g) & [2;6) & [6;10) & [10;14) & [14;18) & [18;22) & [22;26) \\
					\hline Giá trị đại diện & 4 & 8 & 12 & 16 & 20 & 24 \\
					\hline
					Tần số & 2 & 6 & 10 & 13 & 16 & 13  \\ \hline
				\end{tabular}
			\end{center}
			Từ bảng trên, ta tính được số trung bình của mẫu số liệu trên là $\overline{x} \approx 16{,}93$.
			\itemch \textbf{Sai}.\\
			Mẫu số liệu trên có $60$ số liệu, do đó nhóm chứa trung vị của mẫu số liệu là nhóm chứa số liệu ở vị trí thứ 30 và 31 của mẫu số liệu, tức là nhóm $[14;18)$.
			\itemch \textbf{Sai}.\\
			Nhóm tứ phân vị thứ ba là nhóm $[18;22)$. Ta có
			$$Q_3=18+\dfrac{\dfrac{3n}{4}-(2+6+10+13)}{16}\cdot(22-18)=19{,}75.$$
			\itemch \textbf{Sai}.\\
			Độ lệch chuẩn của mẫu số liệu trên là $s\approx5{,}58$.
		\end{itemchoice}
	}
\end{ex}
\begin{ex}%[2D1H5-5]%[2-TN-DS-TLN-THPT-ThienPhuoc-DongNai-HKI-NH24-25 - Phúc Chương]
	Cho hàm số $y=f(x)=\dfrac{ax^2+bx+c}{mx+n}$ có đồ thị như hình vẽ:
	\begin{center}
		\begin{tikzpicture}[scale=.6, font=\footnotesize, line join=round, line cap=round, >=stealth]
			\fill (0,0)node[above left]{$O$} circle (1pt);
			\draw[->] (-3,0)--(7,0)node[above]{$x$};
			\draw[->] (0,-4)--(0,7)node[right]{$y$};
			\begin{scope}
				\clip (-3,-4) rectangle (7,7);
				\draw[smooth,samples=100,domain=-3:1.99] 
				plot(\x,{\x-1+1/(\x-2)});
				\draw[smooth,samples=100,domain=2.01:7] 
				plot(\x,{\x-1+1/(\x-2)});
				\draw[dashed,smooth,samples=100,domain=-3:7] 
				plot(\x,{\x-1});
			\end{scope}
			\draw[dashed] (2,-4)--(2,7)
			(3,0)node[below]{$3$}|-(0,3)node[left]{$3$} (1,0)node[above]{$1$}|-(0,-1)node[left]{$-1$};
		\end{tikzpicture}
	\end{center}
	\choiceTF
{\True $\max\limits_{x\in(-\infty;2)}f(x)=f(1)$}
{\True Hàm số đồng biến trên khoảng $(-\infty;1)$}
{Khoảng cách giữa $2$ điểm cực trị của đồ thị hàm số $y=f(x)$ là $5\sqrt{2}$}
{\True Đồ thị hàm số $y=f(x)$ có đường tiệm cận xiên là $y=x-1$}
	\loigiai{
		\begin{itemchoice}
			\itemch \textbf{Đúng}.\\
			Dựa vào đồ thị của hàm số $y=f(x)$, ta thấy hàm số $y=f(x)$ đồng biến trên $(-\infty;1)$.
			\itemch \textbf{Sai}.\\
			Dựa vào đồ thị, ta thấy đồ thị hàm số có điểm cực tiểu là $A(3;3)$ và điểm cực đại là $B(1;-1)$.\\
			Ta có $AB=\sqrt{(3-1)^2+(3-(-1))^2}=\sqrt{20}=2\sqrt{5}$.
			\itemch \textbf{Đúng}.\\
			Dựa vào đồ thị, ta thấy giá trị lớn nhất của hàm số trên khoảng $(-\infty;2)$ là $-1$ đạt tại $x=1$ hay $\max\limits_{x\in(-\infty;2)}f(x)=f(1)$.
			\itemch \textbf{Đúng.}\\
			Ta thấy, tiệm cận xiên của đồ thị hàm số là đường thẳng đi qua hai điểm có toạ độ $(0;-1)$ và $(1;0)$. Do đó đồ thị hàm số có tiệm cận xiên là $y=x-1$.
		\end{itemchoice}
	}
\end{ex}
\begin{ex}%[2H2V2-6]%[2-TN-DS-TLN-THPT-ThienPhuoc-DongNai-HKI-NH24-25 - Phúc Chương]
	Với hệ trục tọa độ $Oxyz$ (đơn vị: mét) sao cho $O$ nằm trên mặt nước, mặt phẳng $(Oxy)$ trùng với mặt nước, trục $Oz$ hướng lên trên, mội con chim bói cả đang ở vị trí cách mặt nước $2$~m, cách mặt phẳng $(Oxz)$, $(Oyz)$ lần lượt $3$~m và $1$~m phóng thẳng xuống vị trí con cá, biết con cá cách mặt nước $50$~cm, cách mặt phẳng $(Oxz)$, $(Oyz)$ lần lượt là $1$~m và $1{,}5$~m. Gọi $A$ là vị trí con cá, $C$ là vị trí con chim bói cá, điểm $B$ là điểm lúc chim bói cá vừa tiếp xức với mặt nước (xem hình vẽ dưới).
	\begin{center}
		\begin{tikzpicture}[line join=round, line cap=round,scale=1,transform shape, >=stealth]
			\definecolor{columbiablue}{rgb}{0.61, 0.87, 1.0}%màu nước
			\definecolor{arsenic}{rgb}{0.23, 0.27, 0.29}%màu mỏ
			\definecolor{antiquewhite}{rgb}{0.98, 0.92, 0.84}%màu trắng
			\definecolor{cadmiumorange}{rgb}{0.93, 0.53, 0.18}%lông cam
			\definecolor{coolblack}{rgb}{0.0, 0.18, 0.39}%cánh đậm
			\definecolor{brandeisblue}{rgb}{0.0, 0.44, 1.0}%màu xanh đầu
			\definecolor{darkcoral}{rgb}{0.8, 0.36, 0.27}%màu chân
			
			%---------màu vẽ cá
			\definecolor{amber}{rgb}{1.0, 0.49, 0.0}
			% \clip (-3,-3.5) rectangle (3.5,3);
			
			\tikzset{san/.pic={ 
					\path
					(-1.3,-1.5) coordinate (O)
					($(O)+(-142:2)$) coordinate (y)
					($(O)+(0:4.7)$) coordinate (x)
					($(O)+(90:3)$) coordinate (z)
					($(x)+(y)-(O)$) coordinate (t)					
					(-.8,-2.2)coordinate (A)
					(1.6,0.8) coordinate (C)
					($(A)!.13!(C)$) coordinate (B)
					;
					\fill[columbiablue] (O)--(x)--(t)--(y)--cycle;
					
					\foreach\p/\g/\t in {x/110/y, y/150/x, z/-30/z}
					{
						\node at (\p) [shift=(\g:2mm)] {$\t$};
					}					
					\foreach \x/\pos in {A/170,B/0,C/-60,O/150} \fill (\x) circle(1pt) node[{shift=(\pos:0.25)}]{$\x$}; 
					
					\draw[->] (O)--(x) ;
					\draw[->] (O)--(y);
					\draw[->] (O)--(z);
					\draw[dashed] (A)--(B);
					\draw (B)--(C);
					%---------nước
					\draw (-1,-1.7)
					..controls +(-120:.5) and +(-160:.5) ..(0,-2)
					(-.9,-1.9)
					..controls +(-70:.2) and +(-160:.2) ..(-.2,-1.9)
					(-.95,-1.8)
					..controls +(70:.5) and +(30:.5) ..(0,-1.9)
					(.1,-1.6)
					..controls +(-20:.2) and +(30:.2) ..(.2,-1.9)
					(-.7,-1.75)
					..controls +(-170:.2) and +(-160:.3) ..(-.5,-1.9)
					;
			}}			
			\path
			(0,0)pic[scale=1]{san}
			;			
			\tikzset{chim_boi_ca/.pic={
					%==============cánh trái
					\draw[fill=coolblack] %(-1,1.1)..controls +(90:.3) and +(170:.3) ..
					(-.55,1.4)
					..controls +(110:.7) and +(100:.3) ..(-.9,1.8)
					..controls +(135:.3) and +(120:.3) ..(-1.2,1.8)
					..controls +(145:.25) and +(110:.25) ..(-1.45,1.7)
					..controls +(165:.15) and +(85:.15) ..(-1.75,1.63)
					..controls +(-165:.1) and +(85:.1) ..(-1.9,1.5)
					..controls +(165:.15) and +(85:.1) ..(-2.1,1.3)
					..controls +(-165:.1) and +(95:.1) ..(-2.2,1.1)
					..controls +(-160:.1) and +(95:.1) ..(-2.35,1)--(-1,1.1)
					;
					%======--------------------
					%Tô lông đầu
					\def\L{
						(2,.84)
						..controls +(170:.2) and +(25:.3) ..(1,.65)
						..controls +(-145:.5) and +(40:.6) ..(.3,.4)
						..controls +(-140:.3) and +(-60:.3) ..(0,.6)
						..controls +(120:.3) and +(-50:.7) ..(-1,1.1)
						..controls +(90:.3) and +(170:.3) ..(-.55,1.4)%1
						..controls +(-40:.2) and +(140:.1) ..(-.25,1.2)
						..controls +(-70:.2) and +(-160:.35) ..(.15,.85)
						..controls +(75:.7) and +(135:.8) ..(1.9,1.3)--(2.1,1)--cycle
						;
					}
					%\draw[red]\L;
					\fill[brandeisblue] \L;
					%==============================
					\draw[fill=antiquewhite] (2,1.05)
					..controls +(165:.2) and +(-35:.2)..(1.7,1.15)
					..controls +(145:.2) and +(65:.2)..(1.2,1.15)
					..controls +(-60:.1) and +(165:.1)..(1.4,1)
					..controls +(-15:.2) and +(-145:.3)..cycle
					;
					\draw[fill=arsenic] (1.6,1.2)
					..controls +(155:.18) and +(55:.15)..(1.23,1.17)
					..controls +(-75:.2) and +(-95:.2)..cycle
					;
					\fill (1.44,1.14) circle(1mm);					
					\fill[cadmiumorange] (2,1.05)
					..controls +(165:.2) and +(-35:.2)..(1.7,1.15)
					..controls +(145:.2) and +(95:.3)..cycle
					;
					\fill[cadmiumorange] (2,.84)
					..controls +(150:.2) and +(-25:.1) ..(1.7,1)
					..controls +(155:.2) and +(-50:.3) ..(1.2,1.08)
					..controls +(130:.2) and +(50:.2) ..(.75,1)
					..controls +(-130:.2) and +(-10:.2) ..(.21,1)
					..controls +(-120:.1) and +(70:.1) ..(.15,.84)
					..controls +(-20:.3) and +(-150:.3) ..(.8,.85)
					..controls +(-10:.3) and +(150:.5) ..(1.7,.85)
					..controls +(-10:.1) and +(150:.1) ..cycle
					;
					%==================
					%viền đen đầu
					\def\X{
						(0,.6)
						..controls +(120:.3) and +(-50:.7) ..(-1,1.1)
						
						(-.55,1.4)%1
						..controls +(-40:.2) and +(140:.1) ..(-.25,1.2)
						..controls +(-70:.2) and +(-160:.35) ..(.15,.85)
						..controls +(75:.7) and +(135:.8) ..(1.9,1.3)
						;
					}
					\draw[black]\X;
					%====================
					%Tô mỏ
					\def\N{
						(1.9,1.3)
						..controls +(-35:.4) and +(140:.3) ..(3.4,.78)
						..controls +(-175:.2) and +(-10:.3) ..(2,.84)
						..controls +(150:.2) and +(-25:.1) ..(1.7,1)
						..controls +(20:.2) and +(175:.1) ..(2,1.05)
						..controls +(-35:.2) and +(155:.1) ..cycle
						;
					}
					%\draw[red]\N;
					\fill[arsenic] \N;
					%Mỏ
					\def\M{
						(1.9,1.3)
						..controls +(-35:.4) and +(140:.3) ..(3.4,.78)
						..controls +(-175:.2) and +(-10:.3) ..(2,.84)
						..controls +(170:.2) and +(25:.3) ..(1,.65)
						;
					}
					\draw[black]\M;
					%==============================					
					%================Cánh phải
					\draw[fill=coolblack]
					%(-2.6,.95)
					%..controls +(-160:.2) and +(145:.3) ..(-1.6,.5)
					%..controls +(-35:.2) and +(145:.3) ..(-1,-.5)
					%..controls +(-35:.2) and +(-145:.2) ..
					(-.6,-.3)%3
					..controls +(-130:.5) and +(-10:.2) ..(-1,-.95)
					..controls +(170:.1) and +(-10:.15) ..(-1.2,-1.02)
					..controls +(170:.1) and +(-40:.15) ..(-1.45,-.95)
					..controls +(160:.1) and +(-80:.15) ..(-1.6,-.8)
					..controls +(140:.1) and +(-70:.15) ..(-1.75,-.65)
					..controls +(140:.1) and +(-70:.15) ..(-1.9,-.5)
					..controls +(160:.1) and +(-70:.15) ..(-2.05,-.35)
					..controls +(150:.1) and +(-70:.15) ..(-2.25,-.2)
					..controls +(-160:.1) and +(-20:.15) ..(-2.5,-.18)
					..controls +(160:.1) and +(-30:.15) ..(-2.75,-.16)
					..controls +(160:.1) and +(-60:.15) ..(-2.95,-.05)
					..controls +(160:.1) and +(-80:.15) ..(-3.2,.05)
					..controls +(160:.1) and +(-90:.15) ..(-3.35,.15)
					..controls +(160:.1) and +(-95:.15) ..(-3.55,.25)
					..controls +(-160:.2) and +(-145:.2) ..(-3.7,.35)
					..controls +(-160:.2) and +(-160:.4) ..(-3.7,.53)
					..controls +(-160:.2) and +(-160:.3) ..(-3.85,.65)
					..controls +(160:.5) and +(-170:.3) ..(-2.6,.95)
					..controls +(0:.5) and +(120:.3) ..(-.6,-.3)%3
					;
					%===================
					%===================
					\fill[brandeisblue]
					(-.8,-1.1)%lông đuôi xanh
					..controls +(-80:.2) and +(140:.2) ..(-.5,-1.5)
					..controls +(-40:.2) and +(140:.4) ..(-.3,-2.5)
					..controls +(135:.7) and +(-85:.6) ..cycle
					;
					\draw (-.3,-2.5)
					..controls +(135:.7) and +(-85:.6) ..(-.8,-1.1)%lông đuôi xanh
					;
					%=============đuôi
					\draw[fill=brandeisblue] (-.45,-2.3)
					..controls +(-95:.2) and +(95:.2) ..(-.4,-2.8)
					--(-.32,-2.8)--(-.25,-2.2)
					(-.25,-2.2)--(-.32,-2.78)--(-.27,-2.78)--(-.16,-2.2)
					;
					%================
					%Tô lông cam
					\def\C{
						(2,.84)
						..controls +(170:.2) and +(25:.3) ..(1,.65)
						..controls +(-145:.5) and +(40:.6) ..(.3,.4)
						..controls +(-140:.3) and +(-60:.3) ..(0,.6)
						..controls +(120:.3) and +(-50:.7) ..(-1,1.1)%2
						..controls +(130:.2) and +(20:.3) ..(-2.6,.95)
						..controls +(-160:.2) and +(145:.3) ..(-1.6,.5)
						..controls +(-35:.2) and +(145:.3) ..(-1,-.5)
						..controls +(-35:.2) and +(-145:.2) ..(-.6,-.3)%3
						..controls +(-135:.2) and +(100:.2) ..(-.8,-1.1)%lông đuôi xanh
						..controls +(-80:.2) and +(140:.2) ..(-.5,-1.5)
						..controls +(-40:.2) and +(140:.4) ..(-.3,-2.5)%đuôi dưới
						..controls +(40:.4) and +(-150:.5) ..(.5,-1.1)%chân
						..controls +(-20:.2) and +(160:.2) ..(.8,-1.15)
						..controls +(150:.1) and +(-70:.1) ..(.65,-.9)
						..controls +(110:.2) and +(-95:.8) ..(1.26,.6)
						..controls +(75:.1) and +(-160:.2) ..cycle
						;
					}					
					\fill[cadmiumorange] \C;
					\draw[black]\C;					
					\draw[black] (-1,1.1)
					..controls +(130:.2) and +(20:.3) ..(-2.6,.95)				
					(-.6,-.3)%3
					..controls +(-135:.2) and +(100:.2) ..(-.8,-1.1)
					;					
					%======================chân
					\draw[fill=darkcoral]
					(.5,-1.1)
					..controls +(-20:.1) and +(170:.1) ..(1.5,-1.1)%chân
					..controls +(-10:.1) and +(-10:.3) ..(1.2,-1.2)
					;
					%móng 2
					\draw (1.47,-1.15)%chân
					..controls +(-10:.1) and +(120:.1) ..(1.63,-1.25)
					;
					%---------------------
					\draw[fill=darkcoral]
					(.7,-1.15)
					..controls +(40:.1) and +(170:.1) ..(1.3,-1.08)%chân
					..controls +(-10:.1) and +(-10:.1) ..(1.2,-1.2)
					..controls +(170:.1) and +(30:.1) ..(1,-1.2)
					;
					%móng 2
					\draw (1.27,-1.15)%chân
					..controls +(-10:.1) and +(120:.1) ..(1.43,-1.25)
					;
					%------------------------
					\draw[fill=darkcoral]
					(.25,-1.1)
					..controls +(-20:.2) and +(-170:.1) ..(.5,-1.1)%chân
					..controls +(-20:.2) and +(160:.2) ..(.8,-1.15)
					..controls +(-20:.2) and +(160:.2) ..(1.1,-1.2)%móng 1
					..controls +(-20:.1) and +(-50:.1) ..(1.03,-1.25)
					..controls +(130:.05) and +(-10:.05) ..(.8,-1.26)
					..controls +(170:.05) and +(10:.05) ..(.5,-1.28)
					..controls +(-150:.1) and +(-160:.15) ..(.3,-1.25)
					..controls +(160:.1) and +(-80:.1) ..(.1,-1.15)
					;
					%móng 1
					\draw (1.1,-1.25)%móng 1
					..controls +(-20:.1) and +(95:.1) ..(1.2,-1.35)
					;
			}}
			%===========Vẽ cá
			\tikzset{ca/.pic={
					%vây
					\def\V{
						(-.35,.74)
						..controls +(120:.12) and +(40:.22) ..(-.7,.72)--cycle
						(-.7,.32)
						..controls +(-170:.1) and +(10:.1) ..(-.95,.3)
						..controls +(60:.1) and +(-140:.1) ..(-.85,.45)--(-.65,.4)--cycle
						(-.3,.32)
						..controls +(-170:.1) and +(10:.1) ..(-.45,.1)
						..controls +(-40:.1) and +(-110:.1) ..(-.1,.37)--cycle
						;
					}
					\fill[amber] \V;
					\draw\V;					
					%-----------------
					\def\C{
						(-1.25,.83)
						..controls +(-45:.2) and +(130:.2) ..(-1,.58)
						..controls +(35:.2) and +(130:.52) ..(.05,.52)
						..controls +(-90:.1) and +(-110:.1) ..(.05,.52)--(.04,.44)
						..controls +(-150:.5) and +(-40:.2) ..(-1,.53)
						..controls +(-140:.1) and +(40:.1) ..(-1.3,.42)
						..controls +(60:.2) and +(-55:.2) ..cycle
						;
					}					
					\fill[amber] \C;
					\draw\C;
					%-----------------
					\def\Cn{
						(-1,.57)
						..controls +(35:.1) and +(130:.4) ..(.01,.52)
						..controls +(-140:.4) and +(-40:.2) ..cycle
						;
					}
					%\draw[ecru!70!black]\Cn;
					\fill[amber!70] \Cn;
					\draw (-.3,.7)
					..controls +(-120:.1) and +(120:.2) ..(-.3,.34);
					
					\draw[fill=white] (-.22,.55) circle (.08);
					\draw[fill=black] (-.22,.55) circle (.048);					
			}}			
			\path (-0.95,-2.5)pic[xscale=-.35,yscale=.35]{ca}
			(1.6,1)pic[xscale=-.15,yscale=.15,rotate=-40]{chim_boi_ca}; 
		\end{tikzpicture}
	\end{center}
	\choiceTF
{Khoảng cách từ con chim bói cá đến con cá là $2{,}5$ mét}
{$A(1{,}5;1;0{,}5)$}
{\True $C(1;3;2)$}
{\True $\overrightarrow{AB}=\dfrac{1}{5}\overrightarrow{AC}$}
	\loigiai{
		\begin{itemchoice}
			\itemch \textbf{Đúng}.\\
			Do chim bói cá cách mặt phẳng $(Oxz)$ và $(Oyz)$ lần lượt $3$ m và $1$ m, đồng thời cách mặt nước $2$ m nên điểm $C$ có toạ độ $(1;3;2)$.
			\itemch \textbf{Sai}.\\
			Do con cá cách mặt phẳng $(Oxz)$ và $(Oyz)$ lần lượt $1$ m và $1{,}5$ m và cách mặt nước $0{,}5$ m nên điểm $A$ có toạ độ $(1{,}5;1;-0{,}5)$.
			\itemch \textbf{Sai}.\\
			Khoảng cách từ chim bói cá đến con cá là độ dài đoạn $AC$.\\
			$AC = \sqrt{(1-1{,}5)^2+(3-1)^2+(2-(-0{,}5))^2} \approx 3{,}24$ m.
			\itemch \textbf{Đúng}.\\
			Ta có $\vec{AC}=(-0{,}5;2;2{,}5)$.
			Vì $B \in (Oxy)$ nên $B$ có toạ độ $\left(x_B;y_B;0\right)$.\\
			Ta có $\vec{BC}=\left(1-x_B;3-y_B;2-0\right)=\left(1-x_B;3-y_B;2\right)$.\\
			Vì $A$, $B$, $C$ thẳng hàng nên $\vec{AC}=k\vec{BC}$ hay $\heva{&k\left(1-x_B\right)=-0{,}5\\&k\left(3-y_B\right)=2\\&2k=2{,}5} \Leftrightarrow \heva{&k=\dfrac{5}{4}\\&x_B=\dfrac{7}{5}\\&y_B=\dfrac{7}{5}.}$\\
			Suy ra $\vec{AB}=\dfrac{5}{4}\vec{BC}$ hay $\vec{AB}=\dfrac{1}{5}\vec{AC}$.
		\end{itemchoice}
	}
\end{ex}
\Closesolutionfile{ans}%%%
\subsection*{PHẦN III. Trả lời ngắn. (6 câu)}
\setcounter{ex}{0}
\Opensolutionfile{ans}[ans-phanIII]
\begin{ex}%[2D1H5-1]%[2-TN-DS-TLN-THPT-ThienPhuoc-DongNai-HKI-NH24-25]%[Vanle Vo]
	Cho hàm số $y=f(x)=ax^3+bx^2+cx+d$ có đồ thị như sau:
\begin{center}
\begin{tikzpicture}[line join=round,line cap=round, font=\footnotesize,scale=0.8,>=stealth]
	\draw[-stealth](-2.5,0)--(2.5,0)node[above]{$x$};
	\draw[-stealth](0,-3.5)--(0,2)node[right]{$y$};		
	\draw[smooth,samples=300,domain=-2.02:2.02] plot(\x,{(\x)^3-3*\x-1});		
	\fill (0,0) circle(1pt)node[below right]{$O$};
	\draw[dashed](-1,0)|-(0,1) (1,0)|-(0,-3);						
	\foreach \x/\g in {1/90,-1/-110}\fill[black] (\x,0) circle (1pt)+(\g:.3)node{$\x$};
	\foreach \x/\g in {1/0,-3/180}\fill[black] (0,\x) circle (1pt)+(\g:.4)node{$\x$};
	\end{tikzpicture}
\end{center}
	Tính giá trị $f(13)$.
	\shortans[4]{2157}
	\loigiai{
Ta có $y=f(x)=ax^3+bx^2+cx+d\Rightarrow y'=3ax^2+2bx+c$.\\
Dựa vào đồ thị ta có
$$\heva{&f(-1)=1\\&f(1)=-3\\&f'(-1)=0\\&f(1)=0}\Leftrightarrow\heva{&-a+b-c+d=1\\&a+b+c+d=-3\\&3a-2b+c=0\\&3a+2b+c=0}\Leftrightarrow\heva{&a=1\\&b=0\\&c=-3\\&d=-1}$$
Do đó đồ thị đã cho là đồ thị của hàm số $y=f(x)=x^3-3x-1$.\\
Vậy $f(13)=13^3-3\cdot13-1=2157$.
}
\end{ex}
\begin{ex}%[2D3H1-3]%[2-TN-DS-TLN-THPT-ThienPhuoc-DongNai-HKI-NH24-25]%[Tran Ngoc Thanh]
	Thống kê số lượt xem mỗi video trong 35 video của một bạn Youtuber mới lập kênh được một tháng:
\begin{center}
		\begin{tabular}{|c|c|c|c|c|c|}
		\hline
		Số lượt xem & {$[200;500)$} & {$[500;500)$} & $[500;1100)$ & $[1100;1400)$ & $[1400;1700)$ \\
		\hline
		Số video & 4 & 7 & 2 & 15 & 7 \\
		\hline
	\end{tabular}
\end{center}
	Tứ phân vị thứ ba của mẫu số liệu ghép nhóm trên bằng bao nhiêu?
	\shortans{1365}
\loigiai{
	Ta có $ Q_{3}=x_{27} \in [1100; 1400)$.\\
	Suy ra $Q_{3}=1100+\dfrac{\dfrac{3\cdot 35}{4}-13}{15}\cdot 300= 1365$.	
	}
	\end{ex}
\begin{ex}%[2D1V5-1]%[2-TN-DS-TLN-THPT-ThienPhuoc-DongNai-HKI-NH24-25]%[Tran Ngoc Thanh]
	Cho hàm số $y=\dfrac{2x+3}{-3x+5}$. Biết đồ thị hàm số cắt 2 trục tọa độ tại $A\ (a; 0)$ và $B(0; b)$. Tính giá trị biểu thức $T=a b$.\\
	\shortans[4]{-0,9}
	\loigiai{
	Đồ thị hàm số giao với trục $Ox$ tại điểm $\left(\dfrac{-3}{2};0\right)$ và giao với trục $Oy$ tại điểm $\left(0;\dfrac{3}{5}\right)$.\\
	Suy ra $T=a b=\dfrac{-3}{2}\cdot\dfrac{3}{5}=\dfrac{-9}{10}$.
	}
\end{ex}
\begin{ex}%[2D1V3-6]%[2-TN-DS-TLN-THPT-ThienPhuoc-DongNai-HKI-NH24-25]%[Tran Ngoc Thanh]
	Cho hình chữ nhật $ABCD$ nội tiếp nửa đường tròn tâm $O$, bán kính $R=10$ (cạnh $AB$ của hình chữ nhật nằm dọc theo đường kính của đường tròn mà hình chữ nhật đó nội tiếp).
			\begin{center}\begin{tikzpicture}[line join = round, line cap = round,>=stealth,font=\footnotesize,scale=1] 
		\def\R{2}
		\coordinate[label = below:$O$] (O) at (0,0); 
		\coordinate (A) at (-\R,0); 
		\coordinate (B) at ($(A)!2!(O)$);
		\coordinate[label = above right:$C$] (C) at (50:\R); 
		\coordinate[label = above left:$D$] (D) at (130:\R);
		\coordinate[label = below:$A$] (AA) at ($(A)!(D)!(B)$); 
		\coordinate[label = below:$B$] (BB) at ($(A)!(C)!(B)$); 
		\draw (A) arc(180:0:\R)--cycle;
		\draw[fill=cyan!20] (BB)--(C)--(D)--(AA)--cycle;
		\foreach \x in {O,A,B,C,D} \fill[black] (\x) circle (1pt); 
	\end{tikzpicture}	
	\end{center}
	Tính diện tích lớn nhất của hình chữ nhật $ABCD$.\\
	\shortans[4]{100}	
	\loigiai{
		Gọi chiều dài $OA=x\Rightarrow AB=2x$ ($0<x<10$).\\
		$\Rightarrow AD=\sqrt{OD^2-OA^2}=\sqrt{100-x^2}$.\\
		Diện tích hình chữ nhật là $S=AB.AD=2x\sqrt{100-x^2}$.\\
		Xét $f(x)=2x\sqrt{100-x^2}$ trên $(0;10)$, ta có\\
		$$f'(x)=- \dfrac{2 x^{2}}{\sqrt{100 - x^{2}}} + 2 \sqrt{100 - x^{2}}.$$
		Suy ra $f'(x)=0 \Rightarrow x=5 \sqrt{2}$.\\
		Ta có $f\left(5 \sqrt{2}\right)=100;\ f\left(0\right)=0;\ f\left(10\right)=0$ .\\
		Vậy $f(x)$ đạt giá trị lớn nhất tại $x=5 \sqrt{2}$.\\	
	Diện tích lớn nhất là $f\left(5 \sqrt{2}\right)=100$. 
	}
\end{ex}
\begin{ex}%[2H2H2-5]%[2-TN-DS-TLN-THPT-ThienPhuoc-DongNai-HKI-NH24-25]%[Vanle Vo]
	Trong không gian $Oxyz$, cho $\overrightarrow{u}=\overrightarrow{i}+3\overrightarrow{j}-2\overrightarrow{k}$ và $\overrightarrow{v}=(2;-1;1)$. Tính độ dài véctơ $\left[\overrightarrow{u},\overrightarrow{v}\right]$ (kết quả làm tròn đến hàng phần trăm).
\shortans[4]{8,66}
	\loigiai{
$\overrightarrow{u}=\overrightarrow{i}+3\overrightarrow{j}-2\overrightarrow{k}=(1;3;-2)$.\\	
Ta có
$$\left[\overrightarrow{u},\overrightarrow{v}\right]=\left(\begin{vmatrix}
		3 & -2\\
		-1 & 1
		\end{vmatrix}
		;\begin{vmatrix}
		-2 & 1\\
		1 & 2
		\end{vmatrix}
		;\begin{vmatrix}
		1 & 3\\
		2 & -1
		\end{vmatrix}\right)=(1;-5;-7).$$
Độ dài véctơ $\left[\overrightarrow{u},\overrightarrow{v}\right]$ là
$$\left|\left[\overrightarrow{u},\overrightarrow{v}\right]\right|=\sqrt{1^2+(-5)^2+(-7)^2}=\sqrt{75}\approx8{,}66.$$
	}
\end{ex}
\begin{ex}%[2H2H2-2]%[2-TN-DS-TLN-THPT-ThienPhuoc-DongNai-HKI-NH24-25]%[Vanle Vo]
	Trong không gian với hệ tọa độ $Oxyz$, cho hai điểm $A(4; 2; 1)$, $B(-2; -1; 4)$. Gọi $M(a; b; c)$ là điểm thỏa mãn $\overrightarrow{AM}=2\overrightarrow{MB}$. Tính $a+b+c$.
	\shortans[4]{3}
	\loigiai{
Ta có
$$\overrightarrow{AM}=2\overrightarrow{MB}\Leftrightarrow\heva{&x_M-x_A=2\left(x_B-x_M\right)\\&y_M-y_A=2\left(y_B-y_M\right)\\&z_M-z_A=2\left(z_B-z_M\right)}\Leftrightarrow\heva{&x_M=\dfrac{2x_B+x_A}{2}=\dfrac{2\cdot(-2)+4}{3}=0\\&y_M=\dfrac{2y_B+y_A}{2}=\dfrac{2\cdot(-1)+2}{3}=0\\&z_M=\dfrac{2z_B+z_A}{2}=\dfrac{2\cdot4+4}{3}=3.}$$
Suy ra $M(0;0;3)$.\\
Vậy $a+b+c=0+0+3=3$.
	}
\end{ex}
\Closesolutionfile{ans}%%%
\hetde\label{\made}
 \newpage
 \newpage
 \begin{center}% 
 \textbf{\large BẢNG ĐÁP ÁN}
\end{center}
\noindent\textbf{A. ĐÁP ÁN PHẦN I}
\inputansbox{10}{ans-phanI}
\noindent\textbf{B. ĐÁP ÁN PHẦN II}
\inputansbox[2]{2}{ans-phanII}
\noindent\textbf{C. ĐÁP ÁN PHẦN III}
\inputansbox[3]{6}{ans-phanIII}
\end{document}
